\section{The Angels in Scripture}\label{sec:scripture}

Scripture names these spirits by what they do. The Hebrew \emph{mal'akh},
the Greek \gk{angelos}, and the Latin \latin{angelus} each mean one who is
sent with a word. John Paul II gives the sense of both terms: ``angelo
(\latin{angelus}) vuole infatti dire `messaggero'. L'ebraico `malak',
usato nell'Antico Testamento, significa più propriamente `delegato' o
`ambasciatore'\,'' (general audience, 30 July 1986). Aquinas gives the same account: ``Angel means
`messenger.'\,'' Every heavenly spirit is an angel in so far as it makes
divine things known, and the lowest order keeps the common name because it
announces to men immediately (\ST{I q.~108 a.~5 ad~1}).

\subsection{The Name of an Office}

The verse on which the Fathers built this teaching is Ps 103[104]:4, ``Who
makest thy angels spirits: and thy ministers a burning fire,'' which
Hebrews applies to the angels in contrast with the Son (Heb 1:7).
Augustine, expounding it, first sets the existence of the angels on faith
and on the scriptural record: though their appearing is hidden from our
eyes, \latin{tamen esse Angelos novimus ex fide, et multis apparuisse
scriptum legimus}. He then distinguishes nature from office:

\begin{quote}
\latin{Spiritus autem Angeli sunt; et cum spiritus sunt, non sunt angeli;
cum mittuntur, fiunt angeli. Angelus enim officii nomen est, non naturae.
Quaeris nomen huius naturae, spiritus est; quaeris officium, angelus est:
ex eo quod est, spiritus est; ex eo quod agit, angelus est.}
(Augustine, \emph{Enarrationes in Psalmos} 103, s.~1, 15)
\end{quote}

\noindent The Catechism quotes this passage as its answer to the question
``Who are they?'' (CCC~329).

Gregory the Great teaches the same distinction in his homily on the lost
drachma, adding the Greek gloss: \latin{Graeca etenim lingua angeli nuntii,
archangeli vero summi nuntii, vocantur. Sciendum quoque quod angelorum
vocabulum, nomen est officii, non naturae}. The holy spirits are always
spirits, but angels only \latin{tunc~\dots\ cum per eos aliqua nuntiantur},
when something is announced through them; and he cites the same psalm
verse (\emph{Homiliae in Evangelia} 34.8). Aquinas reads the verse under
its other aspect: it is the \emph{sed contra} for the angel's entire
incorporeity, ``Who makes His angels spirits'' (\ST{I q.~50 a.~1}). The
Fathers read the office in \emph{angelos} and Aquinas the nature in
\emph{spiritus}; both readings take the verse of the same creatures (see
\S\ref{sec:q50}).

Because the word names an office, Scripture gives it to men who carry
God's word: the priest ``is the angel of the Lord of hosts'' (Mal 2:7);
the Forerunner is ``my angel'' who prepares the way (Mal 3:1; Matt 11:10);
the seven stars in the hand of the Son of man ``are the angels of the
seven churches'' (Apoc 1:20). Yet the Epistle that calls them all
``ministering spirits'' sets the Son above them by ``a more excellent
name'' (Heb 1:14, 4).

\subsection{The Law and the Historical Books}

\paragraph{The patriarchs.} The first angels named in Scripture guard a
closed door: God ``cast out Adam: and placed before the paradise of pleasure
Cherubims, and a flaming sword, turning every way, to keep the way of the
tree of life'' (Gen 3:24). ``The angel of the Lord'' finds Agar by a
fountain in the wilderness and bids her return (Gen 16:7--11), and calls to
her again from heaven over the dying Ismael (Gen 21:17). The Catechism
gathers these episodes into its summary of the angels' presence ``since
creation and throughout the history of salvation'': ``they closed the
earthly paradise; protected Lot; saved Hagar and her child; stayed
Abraham's hand; communicated the law by their ministry~\dots\ and assisted
the prophets'' (CCC~332).

At the vale of Mambre ``there appeared to him three men standing near to
him'' (Gen 18:2); two of them go on to Sodom as ``the two angels'' (Gen
19:1), sent to destroy the city (Gen 19:13, 15--16). Augustine holds that
the three ``who it is not to be doubted were angels'' were the ones in whom
the Lord appeared, ``although some think that one of them was Christ,''
and answers that opinion from Lot's addressing the two as ``Lord'' in the
singular and from Heb 13:2 (\emph{De civitate Dei} XVI.29). The meal
Abraham served raises Aquinas's question whether angels perform the acts
of life in assumed bodies; their eating was ``not a true eating, but
figurative,'' as Raphael says: ``I seemed indeed to eat and to drink with
you but I use an invisible meat'' (Tob 12:19; \ST{I q.~51 a.~3 ad~5}; see
\S\ref{sec:q51}).

At the sacrifice of Isaac ``an angel of the Lord from heaven called to
him'' (Gen 22:11--12) and speaks in the first person of God (Gen
22:15--16); Augustine takes this as the rule for every such passage: when
the angel speaks, the Lord speaks by him (\emph{De Trinitate}
III.11.23--25). Jacob sees in sleep ``a ladder standing upon the earth
\dots\ the angels also of God ascending and descending by it'' (Gen 28:12),
a vision Christ promises fulfilled in himself (John 1:51). Returning from
Laban, Jacob is met by ``the angels of God'' and names the place Mahanaim
(Gen 32:1--2 [32:2--3]); that night he wrestles with ``a man'' whom Osee
calls an angel (Gen 32:24 [32:25]; Osee 12:4 [12:5]); dying, he blesses in
the name of ``the angel that delivereth me from all evils'' (Gen 48:16).

\paragraph{The Exodus and the Law.} At the burning bush the Douay, with
the Vulgate, reads that ``the Lord appeared to him in a flame of fire out
of the midst of a bush'' (Exod 3:2); Stephen says that ``there appeared to
him, in the desert of mount Sina, an angel in a flame of fire in a bush,''
and Augustine cites Stephen to show that God's voice there came by an
angel (Acts 7:30, 35; \emph{De Trin.}\ III.11.24). ``The angel of God, who
went before the camp of Israel'' passes behind it at the Red Sea (Exod
14:19), and the covenant promises him as guide: ``Behold I will send my
angel, who shall go before thee~\dots\ and my name is in him'' (Exod
23:20--21), the verse Aquinas makes the \emph{sed contra} that some angels
are sent in ministry (\ST{I q.~112 a.~1}; see \S\ref{sec:q112}).

The Law gives the angels a place in Israel's worship: ``two cherubims of
beaten gold'' spread their wings over the propitiatory, whence God speaks
(Exod 25:18--22). The New Testament holds that the Law itself came through
angels: Israel ``received the law by the disposition of angels'' (Acts
7:53); it was ``ordained by angels in the hand of a mediator'' (Gal 3:19);
``the word spoken by angels became steadfast'' (Heb 2:2), which Augustine
reads as testimony that under the Old Testament ``not only those visible
things, but also the word itself, was wrought by angels'' (\emph{De
Trin.}\ III.11.22).

\paragraph{The conquest and the judges.} ``An angel of the Lord stood in
the way against Balaam'' with a drawn sword, seen first by the ass (Num
22:22--23, 31). Before Jericho a man with a drawn sword answers Josue, ``I
am prince of the host of the Lord'' (Jos 5:13--15). An angel greets Gedeon
at the winepress (Judg 6:11--12); the angel who announces Samson's birth
refuses his name---``Why askest thou my name, which is wonderful?''---and
ascends in the flame of Manue's sacrifice (Judg 13:3, 18, 20).

\paragraph{The kings.} An angel twice wakes Elias under the juniper tree
with bread and water (3~Kgs [1~Kgs] 19:5--7). At Dothan the Lord opens the
servant's eyes and ``behold, the mountain was full of horses, and chariots
of fire round about Eliseus'' (4~Kgs [2~Kgs] 6:16--17). In one night ``an
angel of the Lord came, and slew in the camp of the Assyrians a hundred and
eighty-five thousand'' (4~Kgs [2~Kgs] 19:35).

\paragraph{Tobias.} The Book of Tobias is Scripture's fullest narrative of
an angel's ministry to one household (Vulgate numbering). The prayers of Tobias and Sara
``were heard in the sight of the glory of the most high God: And the holy
angel of the Lord, Raphael was sent to heal them both'' (Tob 3:24--25). He appears as a beautiful young man, and Tobias does not know ``that he
was an angel of God'' (Tob 5:5--6); asked his lineage, he gives a man's
name, ``I am Azarias the son of the great Ananias'' (Tob 5:18). At the end
he discloses his ministry and rank: ``I offered thy prayer to the Lord'';
``For I am the angel Raphael, one of the seven, who stand before the
Lord'' (Tob 12:12, 15); John Paul II reads the book as showing that the
angels present men's petitions to God (general audience, 30 July 1986). To
the objection that Raphael, once sent, no longer stands before the Lord,
Aquinas answers with Gregory that those sent on our salvation ``can always
assist and see the face of the Father'' (\ST{I q.~112 a.~3}).

\paragraph{The Machabees.} The second book of Machabees records heavenly
horsemen defending the Temple and Israel. Against Heliodorus ``there
appeared to them a horse, with a terrible rider upon him,'' and ``two other
young men, beautiful and strong, bright and glorious'' who scourged him
(2~Mac 3:25--26). In battle five heavenly horsemen conduct the Jews, two covering Machabeus
with their arms (2~Mac 10:29--30); on the march a horseman goes before
them ``in white clothing, with golden armour'' (2~Mac 11:8).

\subsection{Job, the Psalms, and the Prophets}

\paragraph{Job.} In the heavenly court ``the sons of God came to stand
before the Lord, Satan also was present among them'' (Job 1:6; cf.\ 2:1).
Aquinas answers with Gregory that Satan is not said to have
\emph{assisted} but to be ``present among the assistants,'' for though he
has lost beatitude he retains the angelic nature (\ST{I q.~112 a.~3
ad~3}). God asks Job where he was ``when the morning stars praised me
together, and all the sons of God made a joyful melody'' (Job 38:7).

\paragraph{The Psalms.} The Psalter sings the angels as guardians and as
choir. ``The angel of the Lord shall encamp round about them that fear him:
and shall deliver them'' (Ps 33[34]:8). ``For he hath given his angels
charge over thee; to keep thee in all thy ways. In their hands they shall
bear thee up: lest thou dash thy foot against a stone'' (Ps 90[91]:11--12);
the devil quotes these verses to Christ at the pinnacle of the Temple (Matt
4:6), and Aquinas makes the first of them the \emph{sed contra} for the
guardianship of men by angels (\ST{I q.~113 a.~1}; see \S\ref{sec:q113}).
``Bless the Lord, all ye his angels: you that are mighty in strength, and
execute his word, hearkening to the voice of his orders~\dots\ you
ministers of his that do his will'' (Ps 102[103]:20--21).
The Catechism joins this verse to Matt 18:10: the angels are servants and
messengers of God ``with their whole beings'' (CCC~329).
The angels stand at the summit of the cosmic praise (Ps 148:2), and the
psalmist sings ``in the sight of the angels'' (Ps 137[138]:1;
\emph{Directory} 215). Man was made ``a little less than the angels'' (Ps
8:6; cf.\ Heb 2:7); the manna is ``the bread of angels'' (Ps 77[78]:25),
and the plagues of Egypt were sent ``by evil angels'' (Ps 77[78]:49).

\paragraph{Isaias.} In the year of Ozias's death the prophet sees the Lord
enthroned: ``Upon it stood the seraphims: the one had six wings, and the
other had six wings: with two they covered his face, and with two they
covered his feet, and with two they flew. And they cried one to another,
and said: Holy, holy, holy, the Lord God of hosts, all the earth is full of
his glory'' (Isa 6:2--3). ``One of the seraphims flew to me, and in his hand
was a live coal,'' and touching the prophet's lips said, ``thy iniquities
shall be taken away'' (Isa 6:6--7). Aquinas gives two readings of the sending: following Dionysius (\emph{CH}
13), the angel sent was of an inferior order and was called a Seraph ``in
an equivocal sense, because he came to `kindle' the lips of the prophet'';
or a Seraph acted through an inferior angel (\ST{I q.~112 a.~2 ad~2}).
From Isa 6:3 he also takes the \emph{sed contra} that one order contains
many angels (\ST{I q.~108 a.~3}). The lament over the king of Babylon, ``How art
thou fallen from heaven, O Lucifer, who didst rise in the morning?~\dots\ I
will be like the most High'' (Isa 14:12--14), Aquinas reads ``of the devil
under the figure of the prince of Babylon'' (\ST{I q.~63 a.~5}).

\paragraph{Ezechiel.} By the river Chobar Ezechiel sees ``the likeness of
four living creatures,'' each with four faces and four wings (Ezek 1:5--6),
and later identifies them: ``I understood that they were cherubims'' (Ezek
10:20). The oracle against the king of
Tyre, ``Thou wast the seal of resemblance, full of wisdom, and perfect in
beauty. Thou wast in the pleasures of the paradise of God: every precious
stone was thy covering'' (Ezek 28:12--13), Gregory reads of the first angel
created, whose nine precious stones signify the nine orders with which he
was adorned (\emph{Hom.\ in Ev.}\ 34.7). Aquinas cites the same chapter as
``said to the devil in the person of the King of Tyre'' (\ST{I q.~63
a.~5}) and discusses its ``Thou a cherub stretched out'' (Ezek 28:14) in
asking whether the highest of the fallen was the highest of all (\ST{I
q.~63 a.~7}; see \S\ref{sec:q63}).

\paragraph{Daniel.} ``The angel of the Lord went down with Azarias and his
companions into the furnace'' and made its midst ``like the blowing of a
wind bringing dew'' (Dan 3:49--50, in the Greek portion of Daniel which the
Vulgate carries). Before the Ancient of days ``thousands of thousands
ministered to him, and ten thousand times a hundred thousand stood before
him'' (Dan 7:10). Aquinas makes the verse the \emph{sed contra} for the
great number of the angels (\ST{I q.~50 a.~3}), and he reports two readings
of its proportion: Gregory takes ``thousands of thousands'' in a partitive
sense and holds that those who minister outnumber those who assist, while
Dionysius holds that ``the multitude of angels surpasses all the multitude
of material things'' (\ST{I q.~112 a.~4 ad~2}). Daniel is the first book to
name an angel. ``Gabriel, make this man to understand the vision'' (Dan 8:16); he comes
``flying swiftly'' at the evening sacrifice (Dan 9:21). The angel who
answers Daniel's prayer reports that ``the prince of the kingdom of the
Persians resisted me one and twenty days'' until Michael came to help, and
that he will return to fight that prince with ``Michael your prince''
(Dan 10:13, 20--21). At the end ``shall Michael rise up, the great prince,
who standeth for the children of thy people'' (Dan 12:1). Jerome reads
the prince of Persia as the angel entrusted with that kingdom, who resisted
Israel's release on behalf of his province (\emph{In Danielem} on 10:13;
\S\ref{sec:fl-jerome}), and Gregory as the good angel set over it;
Aquinas follows Gregory and explains the
``resistance'' as the angels' consulting the divine will over the merits
of peoples, ``not that their wills are in opposition'' (\ST{I q.~113
a.~8}). From Dan 10:13 he assigns the guardianship of the human race to
the Principalities, ``or perhaps to the `Archangels' '' (\ST{I q.~113
a.~3}).

\paragraph{Zacharias.} The night visions of Zach 1--6 are mediated
throughout by ``the angel that spoke in me'' (Zach 1:9), who intercedes
for Jerusalem (Zach 1:12).
The high priest Jesus stands ``before the angel of the Lord: and Satan stood
on his right hand to be his adversary. And the Lord said to Satan: The Lord
rebuke thee, O Satan'' (Zach 3:1--2), the rebuke Michael repeats in Jude 9.

\subsection{The Gospels}

\paragraph{The infancy.} The Gospel opens with angelic annunciations. ``An
angel of the Lord, standing on the right side of the altar of incense''
appears to Zachary (Luke 1:11) and names himself: ``I am Gabriel, who
stand before God'' (Luke 1:19). ``The angel Gabriel was sent from God into
a city of Galilee, called Nazareth'' (Luke 1:26); his greeting, ``Hail,
full of grace'' (Luke 1:28), is answered by Mary's ``Behold the handmaid
of the Lord'' (Luke 1:38; see \S\ref{sec:christ-mary}). Gregory explains why
an archangel came: \latin{Ad hoc quippe ministerium summum angelum venire
dignum fuerat, qui summum omnium nunciabat}; it was fitting that the
highest angel should come to this ministry, since he announced the highest
thing of all (\emph{Hom.\ in Ev.}\ 34.8). Joseph is instructed three
times by an angel in sleep (Matt 1:20; 2:13, 19). To the shepherds an
angel stands in the brightness of God, ``and suddenly there was with the
angel a multitude of the heavenly army, praising God'' (Luke 2:9, 13), a
song that ``has not ceased resounding in the Church's praise'' (CCC~333).

\paragraph{The public life.} After the temptation, in which the devil had
quoted the psalm of the guardian angels, ``the devil left him; and behold
angels came and ministered to him'' (Matt 4:11). Christ makes the angels agents of the harvest: ``the reapers are the
angels'' (Matt 13:39); they ``shall gather out of his kingdom all
scandals'' (Matt 13:41) and ``separate the wicked from among the just''
(Matt 13:49). ``There shall be joy before the angels of God upon one
sinner doing penance'' (Luke 15:10): Gregory's Homily 34, preached on
this lection, reads the ten drachmas as the nine orders of angels with man
created tenth, \latin{ut compleretur electorum numerus}, that the number of
the elect might be completed (\emph{Hom.\ in Ev.}\ 34.6), and Aquinas counts that joy among the angels'
accidental reward (\ST{I q.~62 a.~9 ad~3}). Lazarus ``was carried by the
angels into Abraham's bosom'' (Luke 16:22). ``I saw Satan like lightning
falling from heaven'' (Luke 10:18).

Two sayings bear on the angels' nature and their care for men. ``See that
you despise not one of these little ones~\dots\ their angels in heaven
always see the face of my Father'' (Matt 18:10). Aquinas's \emph{sed
contra} for a guardian angel for each man is Jerome's comment on this
verse, ``Great is the dignity of souls, for each one to have an angel
deputed to guard it from its birth''; and he reports Chrysostom's reading
that the text speaks of the highest angels (\ST{I q.~113 a.~2}; a.~3
obj.~1). The risen ``shall neither marry nor be married,'' being ``equal
to the angels'' (Matt 22:30; Luke 20:36), the verse John Paul II takes for
the angels' immortality (general audience, 6 August 1986).

\paragraph{The judgment.} The angels come with the Son of man: he ``shall
send his angels with a trumpet and a great voice'' to gather the elect
(Matt 24:31), though ``of that day and hour no one knoweth: no, not the
angels of heaven'' (Matt 24:36). ``When the Son of man shall come in his
majesty, and all the angels with him'' (Matt 25:31)---the verse with which
the Catechism opens its paragraph on Christ as ``the centre of the angelic
world'' (CCC~331)---he will say to the reprobate, ``Depart from me, you
cursed, into everlasting fire, which was prepared for the devil and his
angels'' (Matt 25:41); Aquinas draws from the verse that all the demons
are subject to the first of them (\ST{I q.~63 a.~8}).

\paragraph{The Passion and Resurrection.} In the garden Christ could ask
the Father for ``more than twelve legions of angels'' (Matt 26:53), and
``there appeared to him an angel from heaven, strengthening him'' (Luke
22:43). At the tomb an angel rolls back the stone, ``his countenance
\dots\ as lightning and his raiment as snow,'' and tells the women, ``He
is not here. For he is risen'' (Matt 28:2--3, 6); Mary Magdalen ``saw two
angels in white, sitting'' where the body had lain (John 20:12). The
Catechism sums up this ministry: the angels ``protect Jesus in his
infancy, serve him in the desert, strengthen him in his agony in the
garden,'' and ``evangelize'' by proclaiming his Incarnation and
Resurrection (CCC~333).

\subsection{Acts and the Apostle}

\paragraph{The Acts of the Apostles.} The Church's first years are
guarded by angels. An angel opens the prison for the apostles (Acts 5:19);
Cornelius sees ``an angel of God coming in unto him'' (Acts 10:3); an
angel frees Peter, ``and the chains fell off from his hands'' (Acts 12:7),
and the disciples, told he is at the gate, answer, ``It is his angel''
(Acts 12:15). On the ship to Rome ``an angel of God, whose I am and whom I
serve, stood by me this night'' (Acts 27:23).

\paragraph{The Pauline letters.} No creature, ``neither death, nor life,
nor angels, nor principalities, nor powers,'' separates us from God's love
(Rom 8:38). The apostles are ``a spectacle to the world and to angels and
to men'' (1~Cor 4:9); the saints ``shall judge angels'' (1~Cor 6:3); the
woman is veiled ``because of the angels'' (1~Cor 11:10); the tongues ``of
men and of angels'' are nothing without charity (1~Cor 13:1); an angel
preaching another gospel is anathema (Gal 1:8). Christ is set ``above all
principality and power and virtue and dominion'' (Eph 1:21), and through
the Church God's wisdom is made known ``to the principalities and powers
in heavenly places'' (Eph 3:10); the same names serve for the fallen,
against whom is our wrestling (Eph 6:12). All were made in the Son,
``whether thrones, or dominations, or principalities, or powers'' (Col
1:16). Against a false cult the Apostle warns, ``Let no man seduce you,
willing in humility and religion of angels'' (Col 2:18; see
\S\ref{sec:devotion}). The Lord will come ``with the voice of an archangel''
(1~Thess 4:15 [4:16]),\footnote{The Douay numbers 1~Thess 4 in seventeen
verses, the Clementine Vulgate and modern editions in eighteen; the
archangel's voice is Douay 4:15.} ``revealed from heaven with the angels
of his power'' (2~Thess 1:7). Timothy is charged ``before God and Christ
Jesus and the elect angels'' (1~Tim 5:21).

\paragraph{Hebrews.} The Epistle to the Hebrews gives Scripture's most
sustained comparison of Christ with the angels. He is ``made so much
better than the angels as he hath inherited a more excellent name than
they'' (Heb 1:4); the first-begotten is adored by them (1:6); ``Are they
not all ministering spirits, sent to minister for them who shall receive
the inheritance of salvation?'' (1:14). The Word ``nowhere doth~\dots\
take hold of the angels: but of the seed of Abraham'' (2:16). The faithful
are come ``to the company of many thousands of angels'' (12:22) and are
bidden to hospitality, ``for by this some, being not aware of it, have
entertained angels'' (13:2). Aquinas takes Heb 1:14 as the principal
objection that all angels are sent, and answers by distinguishing the
assisting and ministering orders (\ST{I q.~112 aa.~2--4}; see
\S\ref{sec:q112}).

\paragraph{The Catholic Epistles.} The prophets ministered the Gospel
``on whom the angels desire to look'' (1~Pet 1:12; \ST{I q.~58 a.~1
ad~2}); Christ is at God's right hand, ``the angels and powers and virtues
being made subject to him'' (1~Pet 3:22). ``God spared not the angels that sinned, but delivered them, drawn down
by infernal ropes to the lower hell'' (2~Pet 2:4); ``the angels who kept
not their principality but forsook their own habitation, he hath reserved
under darkness in everlasting chains'' (Jude 6). The good angels,
``greater in strength and power,'' ``bring not~\dots\ a railing judgment''
(2~Pet 2:11); Jude gives the example: ``When Michael the archangel,
disputing with the devil, contended about the body of Moses, he durst not
bring against him the judgment of railing speech, but said: The Lord
command thee'' (Jude 9). See \S\ref{sec:fall}.

\subsection{The Apocalypse}

The last book shows the angels in the liturgy of heaven and in the
government of history. John hears ``the voice of many angels round about the throne~\dots\ (and
the number of them was thousands of thousands)'' (Apoc 5:11). Four angels
hold the four winds while the servants of God are signed in their
foreheads (7:1--3); seven angels stand in God's presence (8:2), and
another offers a golden censer from which ``the smoke of the incense of
the prayers of the saints ascended up before God'' (8:3--4), the text
under which the Church entrusts the faithful's prayers to the angels'
ministry (Apoc 5:8; \emph{Directory} 215).

``There was a great battle in heaven: Michael and his angels fought with the
dragon, and the dragon fought, and his angels. And they prevailed not:
neither was their place found any more in heaven. And that great dragon was
cast out, that old serpent, who is called the devil and Satan, who
seduceth the whole world. And he was cast unto the earth: and his angels
were thrown down with him'' (Apoc 12:7--9). Aquinas takes the dragon's tail
that ``drew the third part of the stars of heaven'' (12:4) as the
\emph{sed contra} that the sin of the highest angel was the cause of the
others' sinning (\ST{I q.~63 a.~8}). Gregory reads the battle of the end, when the enemy who desired likeness
to God, overthrown by Michael, learns \latin{quia ad Dei similitudinem per
superbiam nullus exsurgat}, that no one rises to God's likeness by pride
(\emph{Hom.\ in Ev.}\ 34.9); John Paul II reads the
chapter of the struggle that grows more violent as the end approaches
(general audience, 20 August 1986).

The book ends with the measure of the honour due to angels. Twice John
falls down to adore the angel and is refused: ``See thou do it not~\dots\
Adore God'' (Apoc 22:8--9; cf.\ 19:10). See \S\ref{sec:liturgy} and
\S\ref{sec:devotion}.

\subsection{The Scriptural Names of the Orders}

Scripture uses nine names for the ranks of the angels, and the Fathers
gathered them into nine orders. Gregory argues from Scripture alone:
angels and archangels, he says, are attested on almost every page of
Scripture, \latin{pene omnes sacri eloquii paginae testantur}, cherubim
and seraphim often in the prophets; Paul names four orders to the
Ephesians and, writing to the Colossians, adds the thrones, so that
\latin{quinque sunt ordines, qui specialiter exprimuntur}; with the rest,
\latin{proculdubio novem esse angelorum ordines inveniuntur}, without
doubt nine orders of angels are found (\emph{Hom.\ in Ev.}\ 34.7). Aquinas's \emph{sed contra} for the
propriety of the names is ``the authority of Holy Scripture wherein they are
so named'' (\ST{I q.~108 a.~5}).

\begin{threestable}{Order}{Scriptural locus (\ST{I q.~108 a.~5}, s.c.)}{Gregory's gloss (\emph{Hom.\ in Ev.}\ 34.8--10)}
Seraphim & Isa 6:2, 6 & \latin{ardentes vel incendentes}: joined to God
with no spirit between, they burn with incomparable love\\
Cherubim & Ezek 1; cf.\ 10:15--20 (also Gen 3:24; Exod 25:18--22) &
\latin{plenitudo scientiae}: the nearer their contemplation of God's
brightness, the fuller their knowledge\\
Thrones & Col 1:16 & God sits in them and through them decrees his
judgments\\
Dominations & Eph 1:21 (also Col 1:16) & they surpass the principalities;
\latin{dominari vero est etiam subjectos quosque possidere}\\
Virtues & Eph 1:21 (also 1~Pet 3:22) & through them signs and miracles are
most often done\\
Powers & Eph 1:21 (also Col 1:16; Rom 8:38) & adverse powers are subject
to them and restrained from tempting men as much as they would\\
Principalities & Eph 1:21 (also Col 1:16; Rom 8:38) & they preside over
the good spirits and dispose what is to be done\\
Archangels & Jude 9 (also 1~Thess 4:15 [4:16]) & \latin{summi nuntii}:
they announce the highest things\\
Angels & ``many places of Scripture'' & \latin{nuntii}: they announce the
lesser things\\
\end{threestable}

Aquinas notes that Dionysius and Gregory agree on the grades of all the
orders except the Principalities and the Virtues, and that ``each of these
placings may claim authority from the words of the Apostle'': Eph 1:21
counts the middle orders upward and places the Virtues between the Powers
and the Dominations, as Dionysius does, while Col 1:16 counts downward from
the Thrones (\ST{I q.~108 a.~6}). Aquinas's comparison of the two
orderings is expounded in \S\ref{sec:q108} and tabulated in
Appendices~\ref{app:orders} and~\ref{app:orderings}.

Three angels bear proper names in the canonical Scriptures. Michael is ``one
of the chief princes'' and ``your prince'' (Dan 10:13, 21), ``the great
prince'' (Dan 12:1), ``the archangel'' (Jude 9), and leader of the angels
against the dragon (Apoc 12:7). Gabriel interprets Daniel's visions (Dan
8:16; 9:21) and announces the births of the Forerunner and of Christ (Luke
1:19, 26). Raphael heals and guides the house of Tobias (Tob 3:25; 12:15).
Gregory holds that the names are given for their ministries to us, and
interprets them: \latin{Michael namque, quis ut Deus; Gabriel autem,
fortitudo Dei; Raphael vero dicitur medicina Dei} (\emph{Hom.\ in Ev.}\
34.8--9). The angel of Manue refuses to give a name (Judg 13:18), and
Raphael, while hidden, gives a man's (Tob 5:18). The Church regulates
the use of further names (\S\ref{sec:devotion}; Appendix~\ref{app:names}); the
Scripture index is Appendix~\ref{app:scripture}.
